\subsection{Comparison With and Without Fusion}
In Table~\ref{tab:main_result_single_col}, comparing the best results for single models and 2 model fusion, we observe that the performance with fusion is generally better than either of its constituent single models for all interfaces.
We compare the effectiveness of fusion by comparing the fusion performance with the best single model performance among the upstream models using the same interface. 
We find that the correlation across different interfaces is quite high, i.e. the fusion performance for all interfaces either improves over single models or they all under-perform the single models.
For instance, for Whisper Small and Data2Vec base fusion on Speaker Verification, we see a consistent improvement of $14.74\%$, $12.36\%$, $8.35\%$ for WS., HConv., and CHConv interfaces respectively. 
On the other hand, for Data2Vec Base and WavLM Base+ fusion on Speaker Verification, the improvements are $-1.2\%$, $-4.17\%$, $-9.49\%$ for WS., HConv., and CHConv interfaces respectively. 
These results suggest that selecting suitable upstream models for fusion is crucial and is more important than the specific selection of interface.


% This indicates the effectiveness of model fusion.
% Specifically, the largest improvement over single model is on
% % However, we do not see an evident additional improvement when fusing with 3 models.
% Furthermore, if we compare the performance of each interface with the best result 


\subsection{Comparison between different Interfaces}
We find that in most cases, the proposed HConv. and CHConv. interfaces perform the best, especially for single models. For fusion, we observe the same trend except for Multilingual ASR on ML-SUPERB. For this task, the proposed methods approximately match the baseline methods.
We hypothesize the reason for this could be the language mismatch between the training and testing distributions for the upstream models.
Due to the fact that the SSL models we investigate are trained only on English, applying them to multilingual ASR is a more challenging task.%, and using a simple HConv. might not provide sufficient modeling power to handle the new languages.

For the two baseline methods, we observe that WS. is more stable than GumD.
Particularly, for Speaker Verification tasks, using GumD. is significantly worse than WS.
We hypothesize that the speaker information is spread across several dimensions in each layer.
% Hence, selecting different layers for different feature dimensions loses the information for speaker information.
Hence, selecting different layers for each feature dimension 
compromises the integrity of the speaker information in the representation.
% loses the information for speaker information.
When comparing HConv. and CHConv., we see that CHConv. further improves the performance over HConv.,
especially for more challenging tasks such as Mono/Multi ASR or the fusion between Speech SL and SSL models.
%This indicates the potential to consider HConv as a foundation interface and we might be able to improve upon HConv in the future by some revision.

\subsection{Fusion among Speech SSL and SL Models}
We also study the effect of fusion among SSL models vs fusion between SL and SSL models. 
We observe that fusing Whisper with other Speech SSL models often yields the best results over fusing only the Speech SSL models together. We hypothesize this is due to the fact that the SL and SSL models use very different training objectives, causing them to learn more heterogeneous representations. In contrast, all of the SSL models are trained using some form of masked prediction and therefore may learn more similar representations from model to model.
Also, the proposed interfaces HConv. and CHConv. consistently outperforms baseline interfaces, indicating that our interfaces are better at fusing between Speech SSL and SL models.
However, we find an exception on LibriSpeech.
In fact, Whisper by itself in the single model experiment is much worse than other Speech SSL models.
We hypothesize this is because the SSL models generally include LibriSpeech in their training set, whereas Whisper does not.


\subsection{Scalability for number of upstream models}
We further extend the number of our upstream models toward 3 models to test their scalability.
Surprisingly, we do not see an evident further improvement compared to 2 model fusion regardless of interface choices.
We hypothesize the reason for this could be that the representations from the upstream models are too similar and including more diverse models in the fusion could provide further improvements, but leave this investigation for future work.

\subsection{Scalability for Large Model Fusion}
Finally, we evaluate our fusion methods when scaling up model size.
To do this, we perform fusion with WavLM Large, Data2Vec Large, and Whisper Large (we also use the English version for LibriSpeech).
From Table~\ref{tab:largeModel_result}, we observe that the performance gap between our proposed interface and baseline is larger than that of the base models, which is consistent with the findings in~\cite{shih24_interface}.
Specifically, for large models, we have a gap of $3.61$  between WS. and HConv. for Mono-1h, and $4.97$ for Emotion Recognition.
On the other hand, we see a gap of $1.01$ between WS. and HConv. for Mono-1h, $0.71$ for Emotion Recognition when fusing the base models.
However, this trend is not consistent with Speaker Verification and LibriSpeech ASR.

% Improvement for SV
% Large model: $0.48$
% Base model: $1.06$


